% === Preface ===
\section*{前言}
\addcontentsline{toc}{section}{前言}

\subsection*{为什么写这本书}

Post-training 是当前大语言模型（LLM）和视觉语言模型（VLM）研究中演化最快的领域。
从 2022 年 InstructGPT 确立 SFT$\rightarrow$RLHF pipeline，
到 2025 年 DeepSeek-R1 证明纯 RL 可以从零涌现推理能力，
再到 2026 年 thinking models 和 computer use 成为行业默认范式——
短短四年间，post-training 的核心方法论已经经历了多次根本性的范式转变。

现有的 survey 多采用静态分类法：按方法类型（SFT、RLHF、DPO）或能力维度（推理、安全、多模态）组织内容。
这种组织方式在某一时间切面上是清晰的，但难以回答一个更本质的问题：
\textbf{这些方法为什么会出现？它们解决了什么问题？又引发了哪些新的问题？}
本书的写作动机正是源于此——
我们希望从\textbf{演化}的视角重新审视 post-training 领域，
追踪``问题 $\rightarrow$ 解决方案 $\rightarrow$ 新问题 $\rightarrow$ 改进方案''的发展链条，
并揭示不同研究方向之间的思想流动。

\subsection*{如何阅读本书}

本书将 post-training 领域划分为七条问题驱动的演化线程（thread），
每条线程对应一个核心问题。以下是几种推荐的阅读路径：

\begin{itemize}[leftmargin=2em]
  \item \textbf{线性阅读}（Thread 1 $\rightarrow$ Thread 7）：
    按顺序阅读全部章节，获得 post-training 领域的完整全景。
    章节之间存在逻辑递进关系——SFT 奠定基础，preference optimization 补充偏好信号，
    safety 施加约束，reasoning 扩展能力边界，multimodal 拓展模态，
    agentic 连接真实世界，efficiency 使一切落地。

  \item \textbf{线程跳读}：
    选择任意感兴趣的线程直接阅读，
    利用文中的彩色交叉引用箭头（{\color{thread2}$\rightarrow$} \textit{见 \S X.Y} 格式）跳转到相关线程的对应小节。
    每条线程都是自包含的，但交叉引用会揭示更深层的联系。

  \item \textbf{全局导航}：
    Introduction 中的 Figure~1 提供了七条线程的全景演化图，
    可作为阅读导航地图——先定位感兴趣的区域，再深入对应章节。
\end{itemize}

\noindent
本书使用统一的 box 系统标注关键信息：
\begin{itemize}[leftmargin=2em]
  \item \textbf{\color{red!60}红色框}（Problem Box）：标识推动领域演化的核心问题
  \item \textbf{\color{blue!60}蓝色框}（Solution Box）：标识针对问题的关键解决方案
  \item \textbf{\color{green!60!black}绿色框}（Evolution Box）：标识方法论的演化路径
  \item \textbf{\color{yellow!70!black}黄色框}（Open Problem Box）：标识尚未解决的开放问题
\end{itemize}

\subsection*{本书的局限性}

本书存在以下已知局限，请读者在阅读时留意：

\begin{enumerate}[leftmargin=2em]
  \item \textbf{知识截止}：本书的文献覆盖截至 2026 年初。
    Post-training 领域的演化速度极快，部分讨论可能在出版时已需要更新。

  \item \textbf{方法导向}：本书聚焦于训练方法论的演化，
    较少涉及工程部署（serving、inference optimization）和产品化实践。

  \item \textbf{文献覆盖}：本书以已发表的学术工作和公开的技术报告为主要参考来源。
    工业界的大量内部实践未被充分覆盖，
    这意味着某些方法的真实应用场景可能比文中描述的更为广泛。
\end{enumerate}

\subsection*{致谢}

感谢所有在 post-training 领域做出贡献的研究者——
正是他们的工作使得这一领域的演化如此丰富和精彩。
本书的写作过程中，与同行的讨论和交流提供了许多宝贵的视角和洞见。
